\CWHeader{Лабораторная работа \textnumero 3 \enquote{Токенизация}}

Нужно реализовать процесс разбиения текстов документов на токены, который потом будет использоваться при индексации. Для этого потребуется выработать правила, по которым текст делится на токены. Необходимо описать их в отчёте, указать достоинства и недостатки выбранного метода. Привести примеры токенов, которые были выделены неудачно, объяснить, как можно было бы поправить правила, чтобы исправить найденные проблемы.

В результатах выполнения работы нужно указать следующие статистические данные:

\begin{itemize}
    \item Количество токенов.
    \item Среднюю длину токена.
\end{itemize}

Кроме того, нужно привести время выполнения программы, указать зависимость времени от объёма входных данных. Указать скорость токенизации в расчёте на килобайт входного текста. Является ли эта скорость оптимальной? Как её можно ускорить?


\pagebreak
